\documentclass[11pt,a4paper]{article}
\usepackage[margin=2.4cm]{geometry}
\usepackage{graphicx}
\usepackage{amsmath}
\usepackage{amssymb}
\usepackage[colorlinks=true,linkcolor=blue,citecolor=blue,urlcolor=blue]{hyperref}
\graphicspath{{../ML/results/}{../Tools/Plots/v1/}}

\newcommand{\pt}{\ensuremath{p_{\mathrm{T}}}}
\newcommand{\msd}{\ensuremath{m_{\mathrm{SD}}}}
\newcommand{\effs}{\ensuremath{\varepsilon_{S}}}
\newcommand{\effb}{\ensuremath{\varepsilon_{B}}}
\newcommand{\drqq}{\ensuremath{\Delta R(q,\bar{q}')}}

\title{\bfseries A boosted hadronic-W tagger from CMS NanoAOD\\
  high-level features}
\author{Siun Kim\\[2pt]
  \small Seoul National University, CMS group}
\date{July 2026}

\begin{document}
\maketitle

\begin{abstract}
\noindent
We study the reconstruction and identification of highly Lorentz-boosted
hadronically decaying W bosons ($W\to q\bar{q}'$) in CMS Run~3 simulation
(2022 + 2022EE, NanoAODv13). A generator-level and detector-level study
establishes the boost scale at which the two decay quarks merge into a single
jet: an AK8 jet captures the full decay above $\pt^{W}\!\approx\!220\,$GeV,
and both quarks merge into one AK4 jet above $\approx 445\,$GeV, in agreement
with the $\Delta R \approx 2m_W/\pt$ expectation. Building on the observation
that merged-W jets carry higher charged-particle multiplicity than quark jets,
we train two machine-learned taggers on jet-level high-level features
available in standard NanoAOD (no particle-flow candidates): a gradient-boosted
decision tree and a lightweight attention model (``ParT-lite'') operating on
softdrop-subjet and secondary-vertex tokens. At 50\% signal efficiency the
nominal taggers reject QCD jets $2.1\times$ better than the classic
$\tau_{21}$+mass-window selection, closing most of the gap to the
constituent-based ParticleNet ($1.8\times$ ahead). A mass-decorrelated
ParT-lite variant matches the mass-shape neutrality of the stored
mass-decorrelated ParticleNet XqqVsQCD discriminant while delivering
$2.2\times$ its background rejection. The charged multiplicity is the
second-most important input in all trained taggers, confirming the
track-multiplicity hypothesis and mirroring the ATLAS three-variable
($m$, $D_2$, $n_{\mathrm{trk}}$) tagger design.
\end{abstract}

\section{Introduction}

At LHC energies, W bosons produced with transverse momentum well above their
mass decay into quark pairs whose opening angle,
\begin{equation}
  \Delta R(q,\bar q') \;\approx\; \frac{2\,m_W}{\pt^{W}},
  \label{eq:dr}
\end{equation}
shrinks below the radius of standard jet clustering algorithms: the two
partons are reconstructed as a single ``merged'' jet. Identifying such jets
--- W tagging --- is a cornerstone of searches with heavy resonances, of
$t\bar{t}$ and diboson measurements, and of any analysis with hadronic W
decays at high boost.

Both LHC collaborations have converged on a common toolbox. CMS Run~2/3
analyses use the softdrop-groomed jet mass \msd{} together with either
cut-based substructure variables ($\tau_{21}=\tau_2/\tau_1$, or the
mass-decorrelated $N_2^{\mathrm{DDT}}$ energy-correlation discriminant) or
constituent-based deep networks (DeepAK8, ParticleNet, and most recently
transformer models of the ParT family). ATLAS employs a three-variable tagger
combining the combined jet mass, the $D_2$ energy-correlation ratio, and
--- notably --- the \emph{charged-track multiplicity} $n_{\mathrm{trk}}$.
Constituent-based networks are the most powerful, but they require
per-particle inputs (PF candidates) that are \emph{not stored} in standard
CMS NanoAOD; they are available only through stored inference scores
evaluated centrally at production time.

This note asks a practical question: \emph{how much of the modern tagging
performance can be recovered from the jet-level high-level features (HLFs)
that standard NanoAOD does store}, in a framework (SKNanoAnalyzer) without
access to PF candidates? The study proceeds in two stages: a kinematic study
of the merge transition (Sec.~\ref{sec:kine}), and the construction and
evaluation of trained taggers (Secs.~\ref{sec:tagger}--\ref{sec:results}).

\section{Samples, jets, and labels}
\label{sec:samples}

All samples are Run~3 (13.6\,TeV) central CMS Monte Carlo in NanoAODv13,
eras 2022 and 2022EE:
\begin{itemize}
  \item \textbf{W-jet signal:} $t\bar t$ semileptonic (\texttt{TTtoLNu2Q},
    powheg) and fully hadronic (\texttt{TTto4Q}, powheg), plus
    \texttt{WW}/\texttt{WZ} (pythia) to populate the high-\pt{} tail.
  \item \textbf{Background:} inclusive QCD multijet samples, \pt-binned from
    170 to 3200\,GeV (nine bins per era).
\end{itemize}
AK8 (anti-$k_{\mathrm{T}}$, $R=0.8$) jets are selected with
$\pt > 200\,$GeV, $|\eta| < 2.4$, tight jet identification, and
$20 < \msd < 250\,$GeV. Jets are labeled with generator truth: a
\emph{W jet} contains both quark daughters of a last-copy hadronic W within
$\Delta R < 0.8$ of the jet axis and no additional hard-process b quark in
the cone; jets that also contain the b from the same top decay
($t \to bq\bar q'$ capture, $\langle\msd\rangle \approx 146\,$GeV) are kept
as a separate class and excluded from both signal and background in training
and evaluation. Background jets are taken exclusively from the QCD samples.
The final dataset contains 4.36M W jets and 11.9M QCD jets, split
60/20/20 into training, validation, and test sets.

\section{The merge transition and the track-multiplicity observation}
\label{sec:kine}

Generator- and detector-level scans of the hadronic-W decay confirm the
expectation of Eq.~\eqref{eq:dr} in full simulation
(Table~\ref{tab:merge}). The AK8 capture threshold and the gen-level
$\drqq < 0.8$ crossing agree within a few GeV, and the naive formula
underestimates both thresholds by only $\sim$10\%.

\begin{table}[htb]
  \centering
  \caption{Boost scale of the merge transition: gen-W \pt{} at which each
    fraction crosses 50\% (TTLJ+WW+WZ combined, 2022+2022EE).}
  \label{tab:merge}
  \begin{tabular}{lcc}
    \hline
    quantity & 50\% at & naive $2m_W/\pt$ \\
    \hline
    $\drqq < 0.8$ (gen)            & $\sim$227\,GeV & 200\,GeV \\
    AK8 jet captures both quarks   & $\sim$220\,GeV & 200\,GeV \\
    two AK4 jets overlap ($\Delta R<0.8$) & $\sim$222\,GeV & --- \\
    $\drqq < 0.4$ (gen)            & $\sim$462\,GeV & 400\,GeV \\
    AK4 merged (both quarks, one jet) & $\sim$444\,GeV & 400\,GeV \\
    \hline
  \end{tabular}
\end{table}

The same study revealed the observable that motivates this tagger: at
matched jet \pt{} (200--400\,GeV), merged-W jets carry a higher
charged-particle multiplicity ($\langle N_{\mathrm{ch}}\rangle = 20.7$) than
gluon jets (17.9), single-quark W jets (15.2), or light-quark jets (13.6).
Charged multiplicity separates the two-prong W jet from \emph{quark} jets,
while the (softdrop) mass separates it from \emph{gluon} jets --- the two
handles are complementary, which is precisely the logic of the ATLAS
$n_{\mathrm{trk}}$-based tagger.

\section{Tagger construction}
\label{sec:tagger}

\paragraph{Inputs.} All inputs are per-jet quantities stored in NanoAODv13:
kinematics and masses (\msd, ungroomed mass, ParticleNet-regressed mass);
$N$-subjettiness ratios $\tau_{21},\tau_{32},\tau_{43}$; energy-correlation
ratios $N_2^{\beta=1}, N_3^{\beta=1}$; the lepton-subjet fraction LSF$_3$;
total/charged/neutral constituent multiplicities; charged/neutral
hadronic and electromagnetic and muon energy fractions; the two softdrop
subjets (four-momenta, b-tag scores, $\Delta R$, momentum sharing $z$); and
secondary vertices within the jet cone (count, summed mass, maximum
decay-length significance, summed track count, and the three most-displaced
vertices individually).

\paragraph{Models.} Two architectures are trained on identical data:
\begin{enumerate}
  \item \textbf{BDT}: XGBoost, 27 flat features, depth 6,
    early-stopped on validation AUC.
  \item \textbf{ParT-lite}: a 132k-parameter attention model inspired by the
    Particle Transformer. The token set --- 2 softdrop subjets and up to 3
    secondary vertices --- passes through three self-attention blocks;
    a class token built from the jet-level HLF vector then attends to the
    tokens (class-attention pooling) and feeds the classification head.
    This transplants the ParT architecture onto the objects NanoAOD
    actually provides.
\end{enumerate}

\paragraph{Training and decorrelation.} Both classes are reweighted flat in
\pt{} during training so the spectrum cannot be learned. Each model is
trained twice: a \emph{nominal} variant with all features, and a
\emph{mass-decorrelated} (MD) variant with all jet-mass-scale features
removed (\msd, masses, subjet masses, subjet $\Delta R$) and the QCD sample
additionally reweighted flat in $(\pt,\msd)$, suppressing mass sculpting by
construction.

\paragraph{Benchmarks.} Four references from the same test set:
$\tau_{21}$ alone; $\tau_{21}$ within a $65<\msd<105\,$GeV window;
$N_2^{\mathrm{DDT}}$ (5\% quantile map in $(\rho,\pt)$, the CMS cut-based
MD standard); and the stored constituent-based ParticleNet scores ---
\texttt{WvsQCD} (nominal ceiling) and the mass-decorrelated
\texttt{XqqVsQCD}.

\section{Results}
\label{sec:results}

\begin{figure}[htb]
  \centering
  \includegraphics[width=0.66\textwidth]{roc_inclusive.png}
  \caption{Background rejection ($1/\effb$) versus signal efficiency for all
    taggers, AK8 \pt{} 200--1200\,GeV, test split.}
  \label{fig:roc}
\end{figure}

\begin{table}[htb]
  \centering
  \caption{Inclusive background rejection at three signal-efficiency points,
    and the QCD \msd{} sculpting (Jensen--Shannon divergence between passing
    and inclusive QCD mass shapes) at $\effb = 5\%$.}
  \label{tab:rej}
  \begin{tabular}{lcccc}
    \hline
    tagger & rej.@$\effs{=}0.3$ & rej.@$\effs{=}0.5$ & rej.@$\effs{=}0.7$
      & JSD@5\%\,\effb \\
    \hline
    ParticleNet WvsQCD (ceiling) & 299 & 86 & 30 & 0.38 \\
    \textbf{ParT-lite (nominal)} & 136 & 48 & 19 & 0.50 \\
    \textbf{BDT (nominal)}       & 134 & 46 & 18 & 0.50 \\
    $\tau_{21}$ + mass window    & 68  & 23 & 9  & --- \\
    $\tau_{21}$                  & 40  & 15 & 7  & --- \\
    \textbf{BDT MD}              & 31  & 13 & 6  & 0.046 \\
    \textbf{ParT-lite MD}        & 32  & 12 & 5.5 & \textbf{0.011} \\
    $N_2^{\mathrm{DDT}}$(5\%)    & 17  & 7  & 3  & 0.001 \\
    ParticleNet XqqVsQCD (MD)    & 15  & 5.4 & 2.5 & 0.010 \\
    \hline
  \end{tabular}
\end{table}

Figure~\ref{fig:roc} and Table~\ref{tab:rej} summarize the discrimination
performance; Table~\ref{tab:wp} lists analysis-ready working points at fixed
inclusive mistag rates.

\paragraph{Nominal taggers.} Both trained HLF taggers reject QCD
$2.1\times$ better than the classic $\tau_{21}$+mass-window selection at
$\effs = 0.5$ (48 vs.\ 23). The remaining factor of 1.8 to ParticleNet is
the irreducible price of not seeing jet constituents. ParT-lite and the BDT
tie inclusively, but ParT-lite is ahead in \emph{every} \pt{} bin
(e.g.\ 42 vs.\ 36 at 500--800\,GeV) --- the inclusive tie is a binning
artifact of the steeply falling spectra.

\paragraph{Mass decorrelation.} The MD ParT-lite reduces mass sculpting by a
factor $\sim$45 relative to the nominal taggers (JSD 0.011 vs.\ 0.50 at 5\%
mistag), statistically compatible with the stored MD ParticleNet XqqVsQCD
(0.010), while providing $2.2\times$ its background rejection (12 vs.\ 5.4)
and $1.7\times$ that of $N_2^{\mathrm{DDT}}$. Among our two MD models the
attention architecture is clearly preferred: at equal rejection the BDT MD
sculpts $\sim$4$\times$ more, and its signal efficiency degrades toward high
\pt{} where ParT-lite MD holds flat (Table~\ref{tab:wp}).

\paragraph{Feature importance.} In both BDT variants the charged
multiplicity is the second-most important feature (16\% of total gain),
behind only \msd{} (nominal) or $\tau_{32}$ (MD). The
track-multiplicity hypothesis of Sec.~\ref{sec:kine} is therefore not only
qualitatively confirmed but is quantitatively a leading source of
discrimination --- reproducing, in CMS NanoAOD, the design logic of the
ATLAS $(m, D_2, n_{\mathrm{trk}})$ tagger.

\begin{table}[htb]
  \centering
  \caption{Working points at fixed inclusive QCD mistag rate: score
    threshold, inclusive signal efficiency, and signal efficiency per AK8
    \pt{} bin (GeV).}
  \label{tab:wp}
  \resizebox{\textwidth}{!}{\begin{tabular}{llcccccc}
\hline
tagger & WP ($\varepsilon_B$) & threshold & $\varepsilon_S$ (incl.) & $\varepsilon_S^{200\text{--}300}$ & $\varepsilon_S^{300\text{--}500}$ & $\varepsilon_S^{500\text{--}800}$ & $\varepsilon_S^{800\text{--}1200}$ \\
\hline
ParT-lite & 5\% & 0.554 & 0.687 & 0.702 & 0.642 & 0.669 & 0.632 \\
ParT-lite & 2.5\% & 0.724 & 0.538 & 0.554 & 0.488 & 0.537 & 0.504 \\
ParT-lite & 1\% & 0.853 & 0.356 & 0.369 & 0.313 & 0.372 & 0.323 \\
ParT-lite & 0.5\% & 0.903 & 0.241 & 0.252 & 0.208 & 0.257 & 0.253 \\
\hline
BDT & 5\% & 0.592 & 0.679 & 0.697 & 0.629 & 0.635 & 0.585 \\
BDT & 2.5\% & 0.754 & 0.527 & 0.547 & 0.470 & 0.491 & 0.446 \\
BDT & 1\% & 0.868 & 0.349 & 0.366 & 0.300 & 0.337 & 0.291 \\
BDT & 0.5\% & 0.912 & 0.238 & 0.251 & 0.200 & 0.243 & 0.215 \\
\hline
ParT-lite MD & 5\% & 0.614 & 0.387 & 0.405 & 0.333 & 0.393 & 0.358 \\
ParT-lite MD & 2.5\% & 0.693 & 0.261 & 0.274 & 0.220 & 0.286 & 0.271 \\
ParT-lite MD & 1\% & 0.767 & 0.140 & 0.146 & 0.122 & 0.183 & 0.198 \\
ParT-lite MD & 0.5\% & 0.808 & 0.083 & 0.083 & 0.078 & 0.132 & 0.148 \\
\hline
BDT MD & 5\% & 0.581 & 0.392 & 0.412 & 0.338 & 0.326 & 0.268 \\
BDT MD & 2.5\% & 0.667 & 0.254 & 0.268 & 0.213 & 0.212 & 0.181 \\
BDT MD & 1\% & 0.749 & 0.128 & 0.135 & 0.108 & 0.121 & 0.098 \\
BDT MD & 0.5\% & 0.793 & 0.072 & 0.074 & 0.066 & 0.082 & 0.068 \\
\hline
ParticleNet WvsQCD & 5\% & 0.906 & 0.768 & 0.791 & 0.708 & 0.635 & 0.593 \\
ParticleNet WvsQCD & 2.5\% & 0.958 & 0.643 & 0.669 & 0.575 & 0.519 & 0.503 \\
ParticleNet WvsQCD & 1\% & 0.984 & 0.471 & 0.494 & 0.407 & 0.377 & 0.392 \\
ParticleNet WvsQCD & 0.5\% & 0.991 & 0.356 & 0.375 & 0.304 & 0.293 & 0.303 \\
\hline
\end{tabular}
}
\end{table}

\begin{figure}[htb]
  \centering
  \includegraphics[width=0.49\textwidth]{sculpting.png}\hfill
  \includegraphics[width=0.49\textwidth]{jsd_vs_effb.png}
  \caption{Left: QCD \msd{} shape of jets passing each tagger at
    $\effb = 5\%$, compared to the inclusive shape. Right: sculpting (JSD)
    versus background efficiency. The nominal taggers carve out the W mass
    peak; the MD variants leave the spectrum nearly unchanged.}
  \label{fig:sculpt}
\end{figure}

\begin{figure}[htb]
  \centering
  \includegraphics[width=0.6\textwidth]{eff_vs_pt.png}
  \caption{Signal efficiency versus AK8 jet \pt{} at the 1\%-mistag working
    point for the nominal taggers and baselines.}
  \label{fig:effpt}
\end{figure}

\section{Summary}

Starting from a first-principles study of the W-boson merge transition
($\pt^{W}\gtrsim 220\,$GeV for AK8, $\gtrsim 445\,$GeV for AK4), we built,
trained, and characterized a family of boosted-W taggers using only
high-level features stored in standard CMS NanoAOD. The trained taggers
double the rejection of the classic $\tau_{21}$+mass selection; a
lightweight attention model over subjet and secondary-vertex tokens matches
a flat BDT inclusively and beats it in every \pt{} bin; and its
mass-decorrelated variant achieves ParticleNet-XqqVsQCD-level mass
neutrality at twice the rejection, making it a competitive analysis-ready
discriminant where the stored MD scores are insufficient. The charged
constituent multiplicity --- the original physics observation that
motivated this work --- proved to be the second-strongest input overall.

The natural next steps are: extension to the 2023(+BPix) eras; a
data/simulation scale-factor measurement in a $t\bar t$-enriched single-lepton
region; and, should PFNano samples become available, a true
constituent-level Particle Transformer to close the remaining factor
$\sim$1.8 to the ParticleNet ceiling.

\vspace{1em}
\noindent\small
Code and results: \texttt{hi/wtag/} (analyzer: SKNanoAnalyzer
\texttt{Wtag.cc}; ML pipeline: \texttt{hi/wtag/ML/}; plans and full result
tables: \texttt{PLAN.md}, \texttt{PLAN\_v2.md}).

\end{document}
